\documentclass[10pt, letterpaper]{article}

% \usepackage{fontspec} 
% \setmainfont{Roboto}

% Packages:
\usepackage[
    ignoreheadfoot, % set margins without considering header and footer
    top=1 cm, % seperation between body and page edge from the top
    bottom=1 cm, % seperation between body and page edge from the bottom
    left=1 cm, % seperation between body and page edge from the left
    right=1 cm, % seperation between body and page edge from the right
    footskip=1.0 cm, % seperation between body and footer
    % showframe % for debugging 
]{geometry} % for adjusting page geometry
\usepackage{titlesec} % for customizing section titles
\usepackage{tabularx} % for making tables with fixed width columns
\usepackage{array} % tabularx requires this
\usepackage[dvipsnames]{xcolor} % for coloring text
\definecolor{primaryColor}{RGB}{0, 79, 144} % define primary color
\usepackage{enumitem} % for customizing lists
\usepackage{fontawesome5} % for using icons
\usepackage{amsmath} % for math
\usepackage[
    pdftitle={Song Lin Chen's CV},
    pdfauthor={Song Lin Chen},
    pdfcreator={LaTeX with RenderCV},
    colorlinks=true,
    urlcolor=primaryColor
]{hyperref} % for links, metadata and bookmarks
\usepackage[pscoord]{eso-pic} % for floating text on the page
\usepackage{calc} % for calculating lengths
\usepackage{bookmark} % for bookmarks
\usepackage{lastpage} % for getting the total number of pages
\usepackage{changepage} % for one column entries (adjustwidth environment)
\usepackage{paracol} % for two and three column entries
\usepackage{ifthen} % for conditional statements
\usepackage{needspace} % for avoiding page brake right after the section title
\usepackage{iftex} % check if engine is pdflatex, xetex or luatex

% Ensure that generate pdf is machine readable/ATS parsable:
\ifPDFTeX
    \input{glyphtounicode}
    \pdfgentounicode=1
    % \usepackage[T1]{fontenc} % this breaks sb2nov
    \usepackage[utf8]{inputenc}
    \usepackage{lmodern}
\fi



% Some settings:
\AtBeginEnvironment{adjustwidth}{\partopsep0pt} % remove space before adjustwidth environment
\pagestyle{empty} % no header or footer
\setcounter{secnumdepth}{0} % no section numbering
\setlength{\parindent}{0pt} % no indentation
\setlength{\topskip}{0pt} % no top skip
\setlength{\columnsep}{0cm} % set column seperation
\makeatletter
\let\ps@customFooterStyle\ps@plain % Copy the plain style to customFooterStyle
\makeatother

\titleformat{\section}{\needspace{4\baselineskip}\bfseries\large}{}{0pt}{}[\vspace{1pt}\titlerule]

\titlespacing{\section}{
    % left space:
    -1pt
}{
    % top space:
    0.3 cm
}{
    % bottom space:
    0.2 cm
} % section title spacing

\renewcommand\labelitemi{$\circ$} % custom bullet points
\newenvironment{highlights}{
    \begin{itemize}[
        topsep=0.10 cm,
        parsep=0.10 cm,
        partopsep=0pt,
        itemsep=0pt,
        leftmargin=0.4 cm + 10pt
    ]
}{
    \end{itemize}
} % new environment for highlights

\newenvironment{highlightsforbulletentries}{
    \begin{itemize}[
        topsep=0.10 cm,
        parsep=0.10 cm,
        partopsep=0pt,
        itemsep=0pt,
        leftmargin=10pt
    ]
}{
    \end{itemize}
} % new environment for highlights for bullet entries
        

\newenvironment{onecolentry}{
    \begin{adjustwidth}{
        0.2 cm + 0.00001 cm
    }{
        0.2 cm + 0.00001 cm
    }
}{
    \end{adjustwidth}
} % new environment for one column entries

\newenvironment{twocolentry}[2][]{
    \onecolentry
    \def\secondColumn{#2}
    \setcolumnwidth{\fill, 4.5 cm}
    \begin{paracol}{2}
}{
    \switchcolumn \raggedleft \secondColumn
    \end{paracol}
    \endonecolentry
} % new environment for two column entries

\newenvironment{header}{
    \setlength{\topsep}{0pt}\par\kern\topsep\centering\linespread{1.5}
}{
    \par\kern\topsep
} % new environment for the header



% save the original href command in a new command:
\let\hrefWithoutArrow\href

% new command for external links:
\renewcommand{\href}[2]{\hrefWithoutArrow{#1}{\ifthenelse{\equal{#2}{}}{ }{#2 }\raisebox{.15ex}{\footnotesize \faExternalLink*}}}


\begin{document}
    \newcommand{\AND}{\unskip
        \cleaders\copy\ANDbox\hskip\wd\ANDbox
        \ignorespaces
    }
    \newsavebox\ANDbox
    \sbox\ANDbox{}

    \begin{header}
        \textbf{\fontsize{24 pt}{24 pt}\selectfont Song Lin Chen}

        
        
        
        \kern 0.25 cm%
        \mbox{\hrefWithoutArrow{https://xsong.us}{\color{black}{\footnotesize\faGlobe}\hspace*{0.13cm}xsong.us}}%
        \kern 0.25 cm%
        
        \mbox{{\color{black}\footnotesize\faMapMarker*}\hspace*{0.13cm}Taiwan}%
        \kern 0.25 cm%
        \AND%
        \kern 0.25 cm%
        \mbox{\hrefWithoutArrow{mailto:ccssll120061203@gmail.com}{\color{black}{\footnotesize\faEnvelope[regular]}\hspace*{0.13cm}ccssll120061203@gmail.com}}%
        \kern 0.25 cm%
        \AND%
        \kern 0.25 cm%
        \mbox{\hrefWithoutArrow{tel:(+886) 905-359-300}{\color{black}{\footnotesize\faPhone*}\hspace*{0.13cm}(+886) 905-359-300}}%
        \kern 0.25 cm%
        \AND%
        \kern 0.25 cm%
        \mbox{\hrefWithoutArrow{https://www.linkedin.com/in/songlinchen/
}{\color{black}{\footnotesize\faLinkedin}\hspace*{0.13cm}songlinchen}}%
        \kern 0.25 cm%
        

    \end{header}

    \vspace{0.3 cm - 0.3 cm}

\section{Experience}

\begin{twocolentry}{\textit{Aug 2025 – Present}}
\textbf{Security-Focused DevOps Engineer} \\
\textit{Pacston Technologies, Inc. — IT Services \& Consulting}
\end{twocolentry}

\begin{onecolentry}
\begin{highlights}

\item Architected secure \textbf{CI/CD pipelines} using \textbf{Bitbucket Pipelines}, \textbf{Docker}, and \textbf{AWS ECR}, enforcing \textbf{immutable image strategy} and reducing manual deployment effort by 80\%.

\item Designed \textbf{AWS cloud infrastructure} (\textbf{IAM}, \textbf{VPC}, \textbf{ALB}, \textbf{ECS Fargate}, \textbf{App Runner}) with \textbf{least-privilege access control} and private subnet segmentation to enhance security posture.

\item Implemented \textbf{Infrastructure as Code} using \textbf{Terraform} and \textbf{Ansible}, integrating automated provisioning into CI/CD workflows and reducing manual infrastructure setup by 75\%.

\item Built \textbf{Python-based automation tooling} for \textbf{release governance}, \textbf{semantic version enforcement}, and PR-driven CHANGELOG generation, reducing release preparation time by 70\%.

\item Developed \textbf{serverless monitoring system} using \textbf{AWS Lambda}, \textbf{DynamoDB}, \textbf{CloudWatch}, and \textbf{SES} for real-time alerting, reducing incident detection time by 50\%.

\end{highlights}
\end{onecolentry}
    
\section{Security Infrastructure \& AI Projects}

\begin{twocolentry}{\textit{Feb 2026}}
\textbf{Vault Secret Management System (Terraform)}
\end{twocolentry}

\begin{onecolentry}
\begin{highlights}
    \item Built secure infrastructure provisioning workflow using Terraform with encrypted Vault secrets, eliminating \texttt{.env} file sharing across 30-person team and 204 services.
    \item Implemented AWS SSM no-key-based authentication, removing 30+ SSH key pairs from server access and reducing unauthorized access surface by 100\%.
    \item Reduced configuration drift risk by enforcing infrastructure automation across 204 services, replacing ad-hoc manual provisioning with auditable, version-controlled workflows.
\end{highlights}
\end{onecolentry}

\begin{twocolentry}{\textit{Dec 2025}}
\textbf{Serverless Security Monitoring Platform (AWS)}
\end{twocolentry}

\begin{onecolentry}
\begin{highlights}
    \item Built serverless real-time monitoring system (AWS Lambda + DynamoDB) covering 200 services, enabling incident detection from \textbf{zero visibility} to fully automated alerting.
    \item Implemented automated anomaly detection with AWS SES and SMS notifications, processing \textbf{\textasciitilde{}60 events/month} and triggering targeted alerts to reduce incident response time by 50\%.
    \item Developed React dashboard to visualize operational metrics across 200 services, replacing manual log inspection with real-time system health visibility.
\end{highlights}
\end{onecolentry}

\begin{twocolentry}{\textit{Oct 2025}}
\textbf{Terraform-Based Infrastructure Platform (AWS)}
\end{twocolentry}

\begin{onecolentry}
\begin{highlights}
    \item Designed modular Terraform architecture provisioning 30 production workloads (20 App Runner, 6 Fargate, 4 Lambda) across VPC, IAM, ECS, and networking resources, reducing manual provisioning effort by 75\%.
\item Implemented least-privilege IAM policies and secure private subnet architecture, eliminating manual access configuration across all 30 services.
\item Integrated Terraform execution into CI/CD workflow, achieving 100\% reproducible infrastructure provisioning and eliminating human error risk.
\end{highlights}
\end{onecolentry}
\begin{twocolentry}{\textit{Sep 2025}}
\textbf{AISecMap – AI-Powered Security Analysis Platform}
\end{twocolentry}

\begin{onecolentry}
\begin{highlights}
\item Architected an AI-driven \textbf{security automation platform} integrating \textbf{RAG}(Ollama), \textbf{static analysis} (Semgrep, Bandit), and \textbf{container scanning} .
\item Built backend services in \textbf{Go} with \textbf{Postgres} and \textbf{Milvus} for vector retrieval and structured threat analysis.
\item Implemented automated vulnerability detection and \textbf{attack path visualization} using React.

\end{highlights}
\end{onecolentry}


    \section{Education}

\begin{twocolentry}{\textit{2023 – 2025}}
\textbf{National Central University} \\
\textit{M.S. in Network Learning Technology}
\end{twocolentry}

\begin{onecolentry}
\begin{highlights}
    \item Participated in \textbf{4 national-level competitions}, including AWS Hackathon and Google Solution Challenge.
    \item Delivered production-style cloud and AI solutions in competitive engineering environments.
\end{highlights}
\end{onecolentry}

\vspace{0.2 cm}

\begin{twocolentry}{\textit{2018 – 2022}}
\textbf{National University of Kaohsiung} \\
\textit{B.S. in Computer Science}
\end{twocolentry}

\begin{onecolentry}
\begin{highlights}
    \item Awarded \textbf{5 technical competition prizes} across software and IoT engineering projects.
    \item Recognized for excellence in system design and full-stack implementation.
\end{highlights}
\end{onecolentry}




\end{document}